设括号串 \(s=s_1s_2\cdots s_n\) ，其中 \(s_i\in\{(\,),)\}\)。定义数列
\[
x_i=\begin{cases}
+1,& s_i=\text{“(”},\\
-1,& s_i=\text{“)”}.
\end{cases}
\]
并令前缀和
\[
S_0:=0,\qquad S_i:=\sum_{k=1}^{i}x_k\quad(1\le i\le n).
\]
记
\[
S_{\min}:=\min_{0\le i\le n}S_i,\qquad S_n=\sum_{k=1}^{n}x_k.
\]
又记左括号与右括号个数分别为
\[
L:=\#\{i:s_i=\text{“(”}\},\qquad R:=\#\{i:s_i=\text{“)”}\},
\]
则 \(L+R=n\) 且 \(S_n=L-R\)。

记 \(\mathrm{LPS}(s)\) 为串 \(s\) 的“最长合法括号子序列”的长度。则有如下等价表述：
\[
\boxed{\;\mathrm{LPS}(s)\;=\;n-S_n+2\,\min\{S_{\min},\,0\}\;}
\]
以及（利用 \(n-S_n=2R\)）
\[
\boxed{\;\mathrm{LPS}(s)\;=\;2\Big(R+\min\{S_{\min},\,0\}\Big)\;.}
\]

若用折线图表示前缀和（将点 \((i,S_i)\) 连线），其最低纵坐标为
\[
Y_{\min}:=S_{\min},
\]
则上式亦可写作
\[
\boxed{\;\mathrm{LPS}(s)\;=\;2\Big(R+\min\{Y_{\min},\,0\}\Big)\;.}
\]
